\documentclass{../../slides-style}
\usetikzlibrary{arrows.meta}

\slidetitle{Коллекции стандартной библиотеки}{17.10.2026}

\begin{document}

    \begin{frame}[plain]
        \titlepage
    \end{frame}

    \section{Коллекции}

    \begin{frame}{Коллекции}
        \framesubtitle{Java Collections Framework}
        \begin{outline}
            \1 Интерфейсы коллекций
            \1 Реализации общего назначения
            \1 Специализированные реализации
                \2 EnumMap с прошлой пары, например
            \1 Многопоточные реализации
            \1 Абстрактные реализации
            \1 Обёртки
            \1 Алгоритмы
            \1 Инфраструктура
        \end{outline}
    \end{frame}

    \begin{frame}{Основные интерфейсы}
        \begin{outline}
            \1 Iterable --- всё, по чему можно итерироваться
            \1 Collection --- просто группа объектов
                \2 Set, SortedSet, NavigableSet --- множества
                    \3 NavigableSet --- расширение SortedSet, добавляющее возможность получить элемент, ближайший к искомому
                \2 List --- упорядоченная коллекция, вводит понятие \enquote{позиция}
                \2 Queue, Deque --- очереди, деки
            \1 Map --- отображение
                \2 SortedMap, NavigableMap
        \end{outline}
    \end{frame}

    \begin{frame}{Реализации общего назначения}
        \begin{footnotesize}
            \begin{tabu} {| X[1 l p] | X[0.7 l p] | X[1 l p] | X[0.7 l p] | X[1 l p] | X[1.3 l p] |}
                \tabucline-
                Интерфейс  & Хеш-таблица  & Массив      & Дерево   & Список      & Хеш-таблица + Список  \\
                \tabucline-
                \everyrow{\tabucline-}
                Set          & HashSet      &             & TreeSet  &             & LinkedHashSet         \\
                List         &              & ArrayList   &          & LinkedList  &                       \\
                Queue, Deque &              & ArrayDeque  &          & LinkedList  &                       \\
                Map          & HashMap      &             & TreeMap  &             & LinkedHashMap         \\
            \end{tabu}

            \vspace{5mm}
            Плюс PriorityQueue --- реализует Queue, но не Deque
        \end{footnotesize}
    \end{frame}

    \begin{frame}{Абстрактные реализации}
        \begin{outline}
            \1 AbstractCollection
                \2 Требуется iterator() с hasNext() и next() и size()
                \2 Для возможности модификации --- переопределить add() и, для итератора, remove()
                \2 Надо конструктор без аргументов и принимающий Collection
                \2 Могут быть переопределены остальные методы, для большей эффективности
            \1 AbstractSet
            \1 AbstractList
            \1 AbstractSequentialList
            \1 AbstractMap
                \2 Требуется реализовать entrySet(), предоставляющий множество пар \mintinline{java}|Map.Entry<K,V>|
                    \3 \enquote{View Collection} --- элементы не хранит, при изменении Map меняется и entrySet
                \2 Для возможности модификации --- переопределить put() и, для итератора, remove()
        \end{outline}
    \end{frame}

    \begin{frame}{Упорядоченные коллекции}
        \begin{outline}
            \1 JEP 431 (Java 21), цель --- упорядочить коллекции с последовательным доступом к элементам
                \2 Раньше были, например, разные методы для доступа к первому и последнему элементам
            \1 \mintinline{java}|SequencedCollection<E>|
                \2 List (и его реализации, например, ArrayList, LinkedList)
                \2 Deque 
            \1 \mintinline{java}|SequencedSet<E>|
                \2 LinkedHashSet, TreeSet
            \1 \mintinline{java}|SequencedMap<K,V>|
                \2 LinkedHashMap
        \end{outline}
    \end{frame}

    \begin{frame}{Многопоточные коллекции}
        \begin{outline}
            \1 Гарантируют корректность работы в многопоточных сценариях
            \1 BlockingQueue --- блокирует вызывающего на put() и take()
                \2 TransferQueue --- расширение BlockingQueue, знающее про потребителей
                    \3 getWaitingConsumerCount(), transfer()
            \1 BlockingDeque
            \1 ConcurrentMap --- обеспечивает fine-grained concurrency (не блокирует вызывающего без крайней нужды)
            \1 ConcurrentNavigableMap
        \end{outline}
    \end{frame}

    \begin{frame}{Инфраструктура}
        \begin{outline}
            \1 Iterator
                \2 hasNext()
                \2 next()
                \2 remove() --- можно не реализовывать
            \1 ListIterator --- двунаправленный, работает с индексами
            \1 Comparable --- задаёт Natural Ordering
            \1 Comparator --- задаёт произвольный порядок
        \end{outline}
    \end{frame}

    \begin{frame}{Диаграмма}
        \begin{center}
            \tikzset{
                box/.style={draw, rounded corners=1.5pt, align=center,
                            inner xsep=2.4pt, inner ysep=2pt, font=\scriptsize},
                gen/.style={draw, {Triangle[length=1.7mm, width=1.6mm, fill=white]}-},
                rea/.style={draw, dashed, {Triangle[length=1.7mm, width=1.6mm, fill=white]}-},
                lbl/.style={anchor=west, font=\scriptsize},
            }
            \resizebox{0.88\textwidth}{!}{%
            \begin{tikzpicture}
                % Ветвь Collection
                \node[box] (iter)    at (4.6,8.35) {{\tiny «interface»}\\ Iterable};
                \node[box] (coll)    at (4.6,7.40) {{\tiny «interface»}\\ Collection};
                \node[box] (seqcoll) at (4.6,6.20) {{\tiny «interface»}\\ Sequenced\\Collection};
                \node[box] (list)    at (4.6,5.00) {{\tiny «interface»}\\ List};
                \node[box] (arrlist) at (4.6,4.00) {ArrayList};
                \node[box] (llist)   at (4.6,3.20) {LinkedList};
                \node[box] (treeset) at (4.6,2.50) {TreeSet};

                % Ветвь Set
                \node[box] (set)      at (1.9,6.45) {{\tiny «interface»}\\ Set};
                \node[box] (seqset)   at (1.9,5.50) {{\tiny «interface»}\\ SequencedSet};
                \node[box] (sortedset) at (1.9,4.55) {{\tiny «interface»}\\ SortedSet};
                \node[box] (navset)   at (1.9,3.60) {{\tiny «interface»}\\ NavigableSet};
                \node[box] (hashset)  at (1.9,2.65) {HashSet};
                \node[box] (lhashset) at (1.9,1.70) {LinkedHashSet};

                % Ветвь Queue
                \node[box] (queue)    at (7.3,6.45) {{\tiny «interface»}\\ Queue};
                \node[box] (deque)    at (7.3,5.50) {{\tiny «interface»}\\ Deque};
                \node[box] (arrdeque) at (7.3,4.55) {ArrayDeque};
                \node[box] (pq)       at (7.3,3.60) {PriorityQueue};

                % Ветвь Map
                \node[box] (map)       at (11.0,7.40) {{\tiny «interface»}\\ Map};
                \node[box] (seqmap)    at (11.0,6.45) {{\tiny «interface»}\\ SequencedMap};
                \node[box] (sortedmap) at (11.0,5.50) {{\tiny «interface»}\\ SortedMap};
                \node[box] (navmap)    at (11.0,4.55) {{\tiny «interface»}\\ NavigableMap};
                \node[box] (treemap)   at (11.0,3.60) {TreeMap};
                \node[box] (hashmap)   at (11.0,2.65) {HashMap};
                \node[box] (lhashmap)  at (11.0,1.70) {LinkedHashMap};

                % Иерархия интерфейсов: от общего (сверху) к частному (снизу)
                \draw[gen] (iter.south) -- (coll.north);
                \draw[gen] (coll.south west) -- (set.north east);
                \draw[gen] (coll.south) -- (seqcoll.north);
                \draw[gen] (coll.east) -- (queue.west);
                \draw[gen] (set.south) -- (seqset.north);
                \draw[gen] (seqset.south) -- (sortedset.north);
                \draw[gen] (sortedset.south) -- (navset.north);
                \draw[gen] (seqcoll.west) -- (seqset.north east);
                \draw[gen] (set.south west) -- (0.75,6.13) -- (0.75,4.55) -- (sortedset.west);
                \draw[rea] (set.west) -- (0.45,6.45) -- (0.45,2.65) -- (hashset.west);
                \draw[gen] (seqcoll.south) -- (list.north);
                \draw[gen] (queue.south) -- (deque.north);
                \draw[gen] (seqcoll.east) -- (deque.north west);
                \draw[gen] (map.south) -- (seqmap.north);
                \draw[gen] (seqmap.south) -- (sortedmap.north);
                \draw[gen] (sortedmap.south) -- (navmap.north);

                % Реализации
                \draw[gen] (hashset.south) -- (lhashset.north);
                \draw[rea] (list.south) -- (arrlist.north);
                \draw[rea] (list.west) -- (3.25,4.50) -- (3.25,3.60) -- (llist.west);
                \draw[rea] (deque.south west) -- (llist.north east);
                \draw[rea] (deque.south) -- (arrdeque.north);
                \draw[rea] (queue.east) -- (9.15,6.45) -- (9.15,3.60) -- (pq.east);
                \draw[rea] (navmap.south) -- (treemap.north);
                \draw[rea] (navset.south east) -- (treeset.north west);
                \draw[gen] (map.east) -- (12.2,7.40) -- (12.2,5.50) -- (sortedmap.east);
                \draw[rea] (map.north east) -- (12.45,7.67) -- (12.45,2.65) -- (hashmap.east);
                \draw[gen] (hashmap.south) -- (lhashmap.north);
            \end{tikzpicture}%
            }
        \end{center}
    \end{frame}

    \begin{frame}{Исключения}
        \begin{outline}
            \1 UnsupportedOperationException --- реализация не поддерживает метод, объявленный в интерфейсе (помечен как optional)
            \1 ConcurrentModificationException --- обнаружено изменение коллекции, когда оно недопустимо 
                \2 например, попытка продолжить обход изменившейся коллекции
                \2 например, модификация коллекции из другого потока
                \2 может не быть брошено --- best effort (надеяться можно, рассчитывать нельзя)
        \end{outline}
    \end{frame}

    \begin{frame}{Утилиты}
        \begin{outline}
            \1 java.util.Collections
                \2 binarySearch, copy, sort, shuffle и т.д.
                \2 emptyList, emptyMap, ...
                \2 singleton, singletonList, ...
                \2 unmodifiableList, unmodifiableMap, ... --- view, а не копия
            \1 java.util.Arrays
                \2 asList --- view в виде списка над массивом, add() делать нельзя
                \2 sort, compare, equals, fill и т.д. для разных примитивных типов
            \1 Фабрики неизменяемых коллекций: List.of, Set.of, Map.of, Map.ofEntries, List.copyOf, Stream.toList
        \end{outline}
    \end{frame}

    \section{Задача}

    \begin{frame}{Задача}
        Переделать список из задачи про хеш-таблицы, чтобы он был генериком и реализовывал стандартный интерфейс \mintinline{java}|List<E>|. Нужны:
        \begin{outline}
            \1 Мутабельность
            \1 Более эффективный итератор, чем по умолчанию
            \1 Тесты
        \end{outline}
    \end{frame}

\end{document}
